%% abtex2-modelo-relatorio-tecnico.tex, v-1.9.2 laurocesar
%% Copyright 2012-2014 by abnTeX2 group at http://abntex2.googlecode.com/ 
%%
%% This work may be distributed and/or modified under the
%% conditions of the LaTeX Project Public License, either version 1.3
%% of this license or (at your option) any later version.
%% The latest version of this license is in
%%   http://www.latex-project.org/lppl.txt
%% and version 1.3 or later is part of all distributions of LaTeX
%% version 2005/12/01 or later.
%%
%% This work has the LPPL maintenance status `maintained'.
%% 
%% The Current Maintainer of this work is the abnTeX2 team, led
%% by Lauro César Araujo. Further information are available on 
%% http://abntex2.googlecode.com/
%%
%% This work consists of the files abntex2-modelo-relatorio-tecnico.tex,
%% abntex2-modelo-include-comandos and abntex2-modelo-references.bib
%%

% ------------------------------------------------------------------------
% ------------------------------------------------------------------------
% abnTeX2: Modelo de Relatório Técnico/Acadêmico em conformidade com 
% ABNT NBR 10719:2011 Informação e documentação - Relatório técnico e/ou
% científico - Apresentação
% ------------------------------------------------------------------------ 
% ------------------------------------------------------------------------

\documentclass[
% -- opções da classe memoir --
12pt,				% tamanho da fonte
openright,			% capítulos começam em pág ímpar (insere página vazia caso preciso)
oneside,			% para impressão em verso e anverso. Oposto a oneside
a4paper,			% tamanho do papel. 
% -- opções da classe abntex2 --
%chapter=TITLE,		% títulos de capítulos convertidos em letras maiúsculas
%section=TITLE,		% títulos de seções convertidos em letras maiúsculas
%subsection=TITLE,	% títulos de subseções convertidos em letras maiúsculas
%subsubsection=TITLE,% títulos de subsubseções convertidos em letras maiúsculas
% -- opções do pacote babel --
english,			% idioma adicional para hifenização
french,				% idioma adicional para hifenização
spanish,			% idioma adicional para hifenização
brazil				% o último idioma é o principal do documento
]{abntex2}

% ---
% Pacotes básicos 
% ---
\usepackage{lmodern}				% Usa a fonte Latin Modern
\usepackage[T1]{fontenc}			% Selecao de codigos de fonte.
\usepackage[utf8]{inputenc}			% Codificacao do documento (conversão automática dos acentos)
\usepackage{lastpage}				% Usado pela Ficha catalográfica
\usepackage{indentfirst}			% Indenta o primeiro parágrafo de cada seção.
\usepackage{color}					% Controle das cores
\usepackage{graphicx}				% Inclusão de gráficos e imagens
\usepackage{microtype} 				% para melhorias de justificação
\usepackage{mathtools}				% Matemática
\usepackage{multirow}				% Tabelas
\usepackage[table,xcdraw]{xcolor}	% Tabelas (cores)

% ---
\usepackage{epigraph}
\usepackage{amssymb}
\usepackage{subcaption}
\usepackage[utf8]{inputenc}
\usepackage[export]{adjustbox}
\usepackage{wrapfig}
\usepackage{amsmath}
\usepackage{pdfpages}
\usepackage{rotating}
\usepackage{pdflscape}



\sloppy
\clubpenalty=10000 % Para evitar linhas orfas
\widowpenalty=10000 % Para evitar linhas viuvas
\hyphenpenalty=10000 % Para não hifenizar

% ---
% Pacotes adicionais, usados apenas no âmbito do Modelo Canônico do abnteX2
% ---
\usepackage{lipsum}				% para geração de dummy text
% ---

% ---
% Pacotes de citações
% ---
\usepackage[brazilian,hyperpageref]{backref}	 % Paginas com as citações na bibl
%\usepackage[alf]{abntex2cite}	% Citações padrão ABNT
\usepackage[alf,abnt-etal-list=0,abnt-etal-cite=6]{abntex2cite}

% --- 
% CONFIGURAÇÕES DE PACOTES
% --- 

% ---
% Configurações do pacote backref
% Usado sem a opção hyperpageref de backref
\renewcommand{\backrefpagesname}{Citado na(s) página(s):~}
% Texto padrão antes do número das páginas
\renewcommand{\backref}{}
% Define os textos da citação
\renewcommand*{\backrefalt}[4]{
	\ifcase #1 %
	Nenhuma citação no texto.%
	\or
	Citado na página #2.%
	\else
	Citado #1 vezes nas páginas #2.%
	\fi}%
% ---

% ---
% Informações de dados para CAPA e FOLHA DE ROSTO
% ---
\titulo{Título da Monografia}
\autor{Nome do Autor}
\local{São Luís}
\data{\the\year}
\orientador{Prof. Dr. Nome do Orientador}

\tipotrabalho{Monografia (Graduação)}
% O preambulo deve conter o tipo do trabalho, o objetivo, 
% o nome da instituição e a área de concentração 
\preambulo{Monografia apresentada ao curso de Ciência da Computação da Universidade Federal do Maranhão, como parte dos requisitos necessários para obtenção do grau de Bacharel em Ciência da Computação.}
% ---


% ---
% Configurações de aparência do PDF final

% alterando o aspecto da cor azul
\definecolor{blue}{RGB}{41, 5, 195}

% informações do PDF
\makeatletter
\hypersetup{
	%pagebackref=true,
	pdftitle={\@title}, 
	pdfauthor={\@author},
	pdfsubject={\imprimirpreambulo},
	pdfcreator={LaTeX with abnTeX2},
	pdfkeywords={abnt}{latex}{abntex}{abntex2}{trabalho acadêmico}, 
	colorlinks=true,       		% false: boxed links; true: colored links
	linkcolor=blue,          	% color of internal links
	filecolor=magenta,      		% color of file links
	urlcolor=blue,
	bookmarksdepth=4
}
\makeatother
% --- 

% --- 
% Espaçamentos entre linhas e parágrafos 
% --- 

% O tamanho do parágrafo é dado por:
\setlength{\parindent}{1.3cm}

% Controle do espaçamento entre um parágrafo e outro:
\setlength{\parskip}{0.2cm}  % tente também \onelineskip

% ---
% compila o indice
% ---
\makeindex
% ---

% ----
% Início do documento
% ----
\begin{document}
	
% Seleciona o idioma do documento (conforme pacotes do babel)
%\selectlanguage{english}
\selectlanguage{brazil}

% Retira espaço extra obsoleto entre as frases.
\frenchspacing 

% ----------------------------------------------------------
% ELEMENTOS PRÉ-TEXTUAIS
% ----------------------------------------------------------
% \pretextual

% ---
% Capa
% ---
\renewcommand{\imprimircapa}{%
\begin{capa}%
\center
\includegraphics[scale=0.25]{logoufma.png} \\ 
\vspace{1cm}
\ABNTEXchapterfont\Large UNIVERSIDADE FEDERAL DO MARANHÃO \\
\ABNTEXsectionfont\Large Curso de Ciência da Computação \\

\vspace*{5cm}

\begin{center}
{\ABNTEXsectionfont\Large\imprimirautor} \\
\vfill
\ABNTEXchapterfont\bfseries\LARGE\imprimirtitulo
\end{center}
\vfill
\large\imprimirlocal \\
\large\imprimirdata
\vspace*{1cm}
\end{capa}
}	

\imprimircapa


% Folha de rosto
% (o * indica que haverá a ficha bibliográfica)
% ---
\imprimirfolhaderosto*
% ---






% ---
% Anverso da folha de rosto:
% ---

%TODO: Quando pegar a ficha catalográfica no SIGAA
%\includepdf{ficha_cat.pdf}

% {

%\ABNTEXchapterfont
%\vspace*{\fill}

%Conforme a ABNT NBR 10719:2011, seção 4.2.1.1.1, o anverso da folha de rosto
%deve conter:
%
%\begin{alineas}
%  \item nome do órgão ou entidade responsável que solicitou ou gerou o
%   relatório; 
%  \item título do projeto, programa ou plano que o relatório está relacionado;
% \item título do relatório;
%  \item subtítulo, se houver, deve ser precedido de dois pontos, evidenciando a
%   sua subordinação ao título. O relatório em vários volumes deve ter um título
%   geral. Além deste, cada volume pode ter um título específico; 
%  \item número do volume, se houver mais de um, deve constar em cada folha de
%   rosto a especificação do respectivo volume, em algarismo arábico; 
%  \item código de identificação, se houver, recomenda-se que seja formado
%   pela sigla da instituição, indicação da categoria do relatório, data,
%   indicação do assunto e número sequencial do relatório na série; 
%  \item classificação de segurança. Todos os órgãos, privados ou públicos, que
%   desenvolvam pesquisa de interesse nacional de conteúdo sigiloso, devem
%    informar a classificação adequada, conforme a legislação em vigor; 
% \item nome do autor ou autor-entidade. O título e a qualificação ou a função
%  do autor podem ser incluídos, pois servem para indicar sua autoridade no
%   assunto. Caso a instituição que solicitou o relatório seja a mesma que o
%   gerou, suprime-se o nome da instituição no campo de autoria; 
%  \item local (cidade) da instituição responsável e/ou solicitante; NOTA: No
%   caso de cidades homônimas, recomenda-se o acréscimo da sigla da unidade da
%   federação.
%  \item ano de publicação, de acordo com o calendário universal (gregoriano),
%  deve ser apresentado em algarismos arábicos.
%\end{alineas}

%\vspace*{\fill}
%}
% ---


%TODO: Após obter as assinaturas da banca, comente as linhas abaixo ou retire
%\includepdf{assinaturas.pdf}


\begin{folhadeaprovacao}

  \begin{center}
    {\ABNTEXchapterfont\large\imprimirautor}

    \vspace*{\fill}\vspace*{\fill}
    \begin{center}
      \ABNTEXchapterfont\bfseries\Large\imprimirtitulo
    \end{center}
    \vspace*{\fill}
    
    \hspace{.45\textwidth}
    \begin{minipage}{.5\textwidth}
        \imprimirpreambulo
    \end{minipage}%
    \vspace*{\fill}
  \end{center}
        
	Trabalho \_\_\_\_\_\_\_\_\_\_ em \imprimirlocal, \today:
	
  \assinatura{\textbf{\imprimirorientador} \\ Orientador} 
  \assinatura{\textbf{Prof. Dr. Primeiro Membro} \\ Examinador}
  \assinatura{\textbf{Profa. Dra. Segundo Membro} \\ Examinador}
%   %\assinatura{\textbf{Professor} \\ Convidado 3}
%   %\assinatura{\textbf{Professor} \\ Convidado 4}
      
  \begin{center}
    \vspace*{0.5cm}
    {\large\imprimirlocal}
    \par
    {\large\imprimirdata}
    \vspace*{1cm}
  \end{center}
  
\end{folhadeaprovacao}


% ---
% Agradecimentos
% ---
\begin{agradecimentos}

Seus agradecimentos.

\end{agradecimentos}
% ---


% Epígrafe
% ---
\newpage
\begin{epigrafe}
    \vspace*{\fill}
	\begin{flushright}
		\textit{"Não fiques em terreno plano. 
		\\Não subas muito alto. 
		\\O mais belo olhar sobre o mundo 
		\\Está a meia encosta."}

Friedrich Nietzsche, em \textit{"A Gaia Ciência" }
	\end{flushright}
\end{epigrafe}
% ---


% ---
% RESUMO
% ---

% resumo na língua vernácula (obrigatório)
%\setlength{\absparsep}{18pt} % ajusta o espaçamento dos parágrafos do resumo
\begin{resumo}
O resumo do trabalho.
	
%\noindent
\textbf{Palavras-chave}: boas palavras chave que identifiquem o trabalho.

\end{resumo}

% resumo em inglês
%	\include{estrutura/resumo_en}
\begin{resumo}[Abstract]
	\begin{otherlanguage*}{english}
	Work abstract.

	\textbf{Keywords}: The best keywords.
		
	\end{otherlanguage*}
\end{resumo}


% ---
\cleardoublepage

% ---
% inserir lista de ilustrações
% ---
\pdfbookmark[0]{\listfigurename}{lof}
\listoffigures*
\cleardoublepage
% ---


% ---
% inserir lista de tabelas
% ---
\pdfbookmark[0]{\listtablename}{lot}
\listoftables*
\cleardoublepage
% ---


% ---
% inserir lista de abreviaturas e siglas
% ---
\begin{siglas}
\item[CC] \textit{Ciência da Computação}
\item[UFMA] \textit{Universidade Federal do Maranhão}
\end{siglas}

% ---

% ---
% inserir lista de símbolos
% ---
%\begin{simbolos}
%  \item[$ \in $] Pertence
%\end{simbolos}
% ---

% ---
% inserir o sumario
% ---

\pdfbookmark[0]{\contentsname}{toc}
\tableofcontents*
\cleardoublepage
	

% ----------------------------------------------------------
% ELEMENTOS TEXTUAIS
% ----------------------------------------------------------
\textual

%\addcontentsline{toc}{chapter}{Resumo}

% ----------------------------------------------------------------------------------------------------------------------------
% Introdução
% ----------------------------------------------------------------------------------------------------------------------------
\chapter{Introdução}
\label{chap:introducao}
Apresente na introdução o problema, a motivação, a necessidade de seu trabalho e os objetivos. Sugere-se incluir trabalhos relacionados e contribuições.

% ----------------------------------------------------------
% ROTEIRO
% ----------------------------------------------------------
%   1. Introdução
%       --> falar do problema, vou te passar os artigos. Pode pegar texto de lá
%       --> objetivos do SILEM
%       --> contribuições

%       2. Fundamentação teórica
%	        --> tudo que foi usado no reconhecimento do analógicos

%       3. Metodologia
%	        --> de maneira simplificada todas as etapas, com os ajustes e parâmetros 
%	        --> explicar aqui a arquitetura do SILEM, tipo, aquela figura que mostra de onde vem, para onde vai

%       4. Resultados
%	        --> o sistema em si e os resultados do reconhecimento do analógicos
%	        --> pequeno estudo de caso de algumas situações

%       5. Conclusões

Algumas dicas no preenchimento da monografia:

Para quebrar um parágrafo, separe eles por uma linha em branco. 

Para referenciar no texto seções ou capítulos, use ref, como em ".... o Capítulo \ref{chap:fundamentacao} apresenta a fundamentação teórica utilizada no trabalho ... ". Lembre-se, os labels dos componentes no texto devem ser únicos!

Para realizar citações, use o comando cite. Quando a citação é indireta, use \cite{associaccao2005abnt}. Quando a citação é direta use \citeonline{associaccao2005abnt}.

Para inserção de figuras, a legenda deve ficar acima, e fonte abaixo. Quando a fonte for o próprio autor, deve-se escrever "Fonte: acervo do autor". Para referencia a figura, use o ref, como " .... Figura \ref{fig:ex1} .... ".

\begin{figure}[ht!]
\centering
\caption{Um exemplo de posicionamento de figura.}
\includegraphics[scale=0.35]{figuras/fig1.png}
\legend{Fonte: \cite{associaccao2005abnt}.}
\label{fig:ex1}
\end{figure}

Uma das formas, em latex, de colocar múltiplas figuras lado a lado é usando o array.

\begin{figure}[ht!]
\centering
\caption{Um exemplo de vetor de figuras.}
$
\begin{array}{cc}
     \includegraphics[scale=0.25]{figuras/fig1.png} & 
     \includegraphics[scale=0.25]{figuras/fig1.png}\\
     (a) &
     (b)
\end{array}
$
\legend{Fonte: acervo do autor.}
\label{fig:ex1}
\end{figure}

Equações numeradas devem ser inseridas usando o comando equation. Caso precise referenciá-las, use ref, com ".... Equação \ref{eq:r_mim_max} ... "

\begin{equation}
\centering
r = (x_{max} - x_{min})G(x) + x_{min} 
\label{eq:r_mim_max}
\end{equation}

Tabelas devem sempre ter a legenda na parte acima. Para referenciá-las, use ref, com ".... Tabela \ref{tab:resultados_1} ... "



\begin{table}[htb]
		\centering
		\caption{Resultado do ....}
		\label{tab:resultados_1}
		\begin{tabular}{cccc}
			\hline
    \textbf{Resultado }& & \textbf{Acertos/Total} & \textbf{Acurácia } \\\hline

	\textbf{Acertou tudo}&    & 245/361 & 67,86\% \\
	\textbf{Errou 1}     &    & 69/361  & 19,11\% \\
	\textbf{Errou 2}     &    & 18/361  & 4,98\%  \\
	\textbf{Errou 3}     &    & 9/361   & 2,49\%  \\
	\textbf{Errou 4}     &    & 5/361   & 1,38\%  \\
	\textbf{Errou tudo}  &    & 15/361  & 4,15\%  \\ \hline
	\end{tabular}
\end{table}

% ----------------------------------------------------------
% Trabalhos relacionados
% ----------------------------------------------------------
\section{Trabalhos Relacionados}
\label{sec:relacionados}

Apresente trabalhos que se assemelham ao proposta. Busque identificar diferenças entre os trabalhos e como o seu trabalho se posiciona em relação a estes.

% ----------------------------------------------------------
% O objetivo do trabalho
% ----------------------------------------------------------
\section{Objetivos}
\label{sec:objetivos}
O objetivo geral desse trabalho é ....

\subsection{Objetivos Específicos}
\label{sec:objetivos_especificos}

Destaca-se como objetivos específicos deste trabalho:

\begin{itemize}
    \item A
    \item B
\end{itemize}

\section{Contribuições}
\label{sec:contrib}

Destaca-se como contribuições deste trabalho:

\begin{itemize}	
       \item A
       \item B
\end{itemize}

\section{Organização do Trabalho}
\label{sec:organizacao}

Este trabalho está organizado em cinco capítulos, de forma a apresentar o conteúdo mais claramente, conforme os parágrafos a seguir.

Capítulo 1 apresenta ....
Capítulo 2 aborda os conteúdos de ....

% ----------------------------------------
% Fim da Introdução
% ----------------------------------------

% ----------------------------------------
% Fundamentação Teórica
% ----------------------------------------
	
\chapter{Fundamentação Teórica}
\label{chap:fundamentacao}
Apresente todo o referencial teórico utilizado como base para a realização do trabalho.


% ---------------------------------------------------
% Fim da Fundamentação
% ---------------------------------------------------

\chapter{Metodologia}
\label{chap:metodologia}
Apresente e explique em detalhes a metodologia proposta. O nome do capítulo pode variar de acordo com o tema de monografia, mas é neste capítulo que é apresentado o que foi feito.

% ---------------------------------------------------
% Fim Metodologia
% ---------------------------------------------------


% --------------------------------------------
% Resultados
% --------------------------------------------
\chapter{Resultados}
\label{chap:resultados}
Apresente os resultados. Analise e discuta sobre os mesmos. Se possível, apresente estudos de caso.

% --------------------------------------------
% Fim Resultados
% --------------------------------------------
	
% --------------------------------------------
% Conclusão
% --------------------------------------------
\chapter{Conclusão}
\label{chap:conclusao}

Aponte na conclusão:
\begin{itemize}
    \item Resumo geral do que foi apresentado no trabalho com um esboço dos resultados atingidos;
    \item Aponte para os objetivos propostos quais foram atingidos e quais não foram;
    \item Determine o quanto o trabalho apresenta de contribuição para a área
    \item Para os objetivos não realizados ou opcionais, aponte trabalhos futuros
\end{itemize}

% --------------------------------------------
% Fim Conclusão
% --------------------------------------------
	
% ---
% ---
\phantompart
	
% ----------------------------------------------------------
% ELEMENTOS PÓS-TEXTUAIS
% ----------------------------------------------------------
\postextual
	
% ----------------------------------------------------------
% Referências bibliográficas
% ----------------------------------------------------------
\bibliography{referencias}
	
% ----------------------------------------------------------
% Glossário
% ----------------------------------------------------------
%
% Consulte o manual da classe abntex2 para orientações sobre o glossário.
%
%\glossary
	
% ----------------------------------------------------------
% Apêndices
% ----------------------------------------------------------
	
% ---
% Inicia os apêndices
% ---
%\begin{apendicesenv}
	
% Imprime uma página indicando o início dos apêndices
%\partapendices
% ----------------------------------------------------------
%\chapter{Quisque libero justo}
% ----------------------------------------------------------
	
%\lipsum[50]
	
% ----------------------------------------------------------
%\chapter{Nullam elementum urna vel imperdiet sodales elit ipsum pharetra ligulaac pretium ante justo a nulla curabitur tristique arcu eu metus}
% ----------------------------------------------------------
%\lipsum[55-57]
	
%\end{apendicesenv}
% ---
	
	
% ----------------------------------------------------------
% Anexos
% ----------------------------------------------------------
	
% ---
% Inicia os anexos
% ---
%	\begin{anexosenv}
%	
%	% Imprime uma página indicando o início dos anexos
%	%\partanexos
%	
%	% ---
%	\chapter{texto}
%
%	\end{anexosenv}
	
\end{document}
